\chapter{Criterio operativo de configuración}

\section{Configuración operativa}
\begin{definition}
Se llama configuración operativa a toda pieza estructural que, dentro del canon intermedio, puede ser tratada como unidad disciplinada de lectura, comparación, transición y articulación.
\end{definition}

\begin{observation}
No toda configuración operativa es todavía un régimen, una familia o una demo. Su función en este nivel es servir como unidad tratable sin perder herencia canónica.
\end{observation}

\section{Condición mínima de configuración}
\begin{definition}
Una pieza alcanza estatus de configuración operativa cuando concurren simultáneamente:
\begin{enumerate}
  \item dependencia visible respecto del núcleo basal integrado en Tomo I;
  \item legibilidad estructural suficiente;
  \item umbral operativo no colapsado;
  \item posibilidad de comparación o distinción disciplinada;
  \item función reconocible dentro del canon intermedio.
\end{enumerate}
\end{definition}

\begin{observation}
Estas condiciones no exigen todavía clasificación exhaustiva ni validación experimental. Exigen únicamente que la pieza ya pueda ser tratada como unidad operativa real.
\end{observation}

\section{Configuración y nivel}
\begin{definition}
Toda configuración operativa debe poder ubicarse en un nivel jerárquico legible o en una transición traducible entre niveles del canon.
\end{definition}

\begin{observation}
Esta cláusula evita que entren al Tomo II piezas estructuralmente indeterminadas o sin lugar canónico reconocible.
\end{observation}

\section{Configuración y comparabilidad}
\begin{definition}
Una configuración operativa es comparable con otra solo si ambas conservan compatibilidad suficiente de herencia, umbral y nivel dentro del dominio declarado.
\end{definition}

\begin{observation}
La comparabilidad entre configuraciones no elimina su distinción operativa. Solo permite tratarlas bajo una relación disciplinada común.
\end{observation}

\section{Configuración y articulación}
\begin{definition}
Una configuración operativa es articulable cuando puede incorporarse a una estructura jerárquica sin perder su función propia ni colapsar en una lista indiferenciada de piezas.
\end{definition}

\begin{observation}
La articulación no debe absorber la configuración hasta volverla irreconocible. Debe conservar su papel dentro del orden canónico.
\end{observation}

\section{Condición mínima de fallo}
\begin{definition}
Se considera fallido el estatus de configuración operativa cuando una pieza pierde herencia visible, umbral suficiente, función reconocible o capacidad de relación disciplinada con otras piezas del mismo canon.
\end{definition}

\begin{observation}
Esta condición protege al Tomo II frente a objetos verbales o híbridos que todavía no alcanzan unidad operativa real.
\end{observation}

\section{Cláusula de alcance}
\begin{observation}
El criterio operativo de configuración fija la unidad mínima tratable del canon intermedio. No constituye todavía una clasificación completa de familias, regímenes o demostraciones.
\end{observation}


